% GENERATED FILE. Edit canonical lessons, then run npm run generate:books.
% canonical-source-hash: fnv1a64:f3fb235973c94c7e
% canonical-lessons: AR-C341-faris, AR-C341-khudha, AR-C341-hudna, AR-C341-muahada, AR-C341-khana

\chapter{Knights and treaties}
\label{ch:ar-knights-and-treaties}
\hlchaptermodality{341}

\begin{chapteropening}
\textbf{By the end of this chapter:} \emph{I can say a knight, a helmet, a truce, a treaty and to betray.}

Complete the last lesson of chapter 341: I can say a knight, a helmet, a truce, a treaty and to betray.
\end{chapteropening}

\section[\texorpdfstring{\ar{فارس} (fāris)}{فارس (fāris)}]{\ar{فارس} (fāris) — a knight, a horseman}

\label{lesson:AR-C341-faris}



\begin{quote}
Before the new one: say the Arabic for a royal court, then the Arabic for to take office.
\end{quote}

\subsection*{You'll want to know: \ar{فارس}}
\textbf{\ar{فارس}} — \emph{fāris} — \textquotedblleft{}a knight, a horseman\textquotedblright{}.

\begin{quote}
\ar{ركب} \ar{الفارس} \ar{حصانه} \ar{الأبيض}.
\end{quote}

\begin{quote}
\emph{rakiba al-fāris ḥiṣānahu al-abyaḍ.}
\end{quote}

\begin{quote}
The knight rode his white horse.
\end{quote}

Plural: \ar{فرسان}, \emph{fursān}.

\ar{فارس} \ar{الأحلام}: Prince Charming.

\begin{cousinweb}
\textbf{\ar{فارس}} (\emph{fāris}) is built on the native Arabic root \emph{f-r-s} (Lane), the root of \emph{faras}, a horse. \emph{Fāris} is the man on the horse; \emph{furūsiyya} is horsemanship and knighthood.
\end{cousinweb}

\begin{grammarlens}[title={What you've built}]
One more word for telling what happened.
\end{grammarlens}

\subsection*{Your turn}
\begin{itemize}
\raggedright
  \item \emph{Say it:} \emph{fāris}
  \item \emph{Say it:} \emph{fāris}, once more
  \item \emph{Say it:} \emph{rakiba al-fāris ḥiṣānahu al-abyaḍ.} — The knight rode his white horse.
  \item \emph{Recall:} say the Arabic for a royal court, then the Arabic for to take office, then say \emph{fāris} again
\end{itemize}

\begin{tcolorbox}[breakable,colback=teal!4,colframe=teal!35!black,title={Before you move on}]
What does \ar{فارس} mean? (\textquotedblleft{}A knight, a horseman\textquotedblright{}.) Say it once more.
\end{tcolorbox}
\section[\texorpdfstring{\ar{خوذة} (khūdha)}{خوذة (khūdha)}]{\ar{خوذة} (khūdha) — a helmet}

\label{lesson:AR-C341-khudha}



\begin{quote}
Before the new one: say the Arabic for to take office, then the Arabic for a knight.
\end{quote}

\subsection*{You'll want to know: \ar{خوذة}}
\textbf{\ar{خوذة}} — \emph{khūdha} — \textquotedblleft{}a helmet\textquotedblright{}.

\begin{quote}
\ar{لبس} \ar{العامل} \ar{خوذة} \ar{صفراء}.
\end{quote}

\begin{quote}
\emph{labisa al-ʿāmil khūdha ṣafrāʾ.}
\end{quote}

\begin{quote}
The worker put on a yellow helmet.
\end{quote}

\begin{cousinweb}
\textbf{\ar{خوذة}} (\emph{khūdha}) is usually traced to Persian \emph{khōd}, \textquotedblleft{}a helmet\textquotedblright{}, an old borrowing into Arabic.
\end{cousinweb}

\begin{grammarlens}[title={What you've built}]
One more word for telling what happened.
\end{grammarlens}

\subsection*{Your turn}
\begin{itemize}
\raggedright
  \item \emph{Say it:} \emph{khūdha}
  \item \emph{Say it:} \emph{khūdha}, once more
  \item \emph{Say it:} \emph{labisa al-ʿāmil khūdha ṣafrāʾ.} — The worker put on a yellow helmet.
  \item \emph{Recall:} say the Arabic for to take office, then the Arabic for a knight, then say \emph{khūdha} again
\end{itemize}

\begin{tcolorbox}[breakable,colback=teal!4,colframe=teal!35!black,title={Before you move on}]
What does \ar{خوذة} mean? (\textquotedblleft{}A helmet\textquotedblright{}.) Say it once more.
\end{tcolorbox}
\section[\texorpdfstring{\ar{هدنة} (hudna)}{هدنة (hudna)}]{\ar{هدنة} (hudna) — a truce}

\label{lesson:AR-C341-hudna}



\begin{quote}
Before the new one: say the Arabic for a knight, then the Arabic for a helmet.
\end{quote}

\subsection*{You'll want to know: \ar{هدنة}}
\textbf{\ar{هدنة}} — \emph{hudna} — \textquotedblleft{}a truce\textquotedblright{}.

\begin{quote}
\ar{وقع} \ar{الطرفان} \ar{هدنة} \ar{لمدة} \ar{شهر}.
\end{quote}

\begin{quote}
\emph{waqqaʿa aṭ-ṭarafān hudna li-muddat shahr.}
\end{quote}

\begin{quote}
The two sides signed a truce for a month.
\end{quote}

\begin{cousinweb}
\textbf{\ar{هدنة}} (\emph{hudna}) is a native Arabic word built on the root \emph{h-d-n}, recorded in \emph{hudna ʿalā dakhan}, an illusory truce (Lane). The root is usually explained as \textquotedblleft{}to be still, to be quiet\textquotedblright{}: a time when the fighting is still.
\end{cousinweb}

\begin{grammarlens}[title={What you've built}]
One more word for telling what happened.
\end{grammarlens}

\subsection*{Your turn}
\begin{itemize}
\raggedright
  \item \emph{Say it:} \emph{hudna}
  \item \emph{Say it:} \emph{hudna}, once more
  \item \emph{Say it:} \emph{waqqaʿa aṭ-ṭarafān hudna li-muddat shahr.} — The two sides signed a truce for a month.
  \item \emph{Recall:} say the Arabic for a knight, then the Arabic for a helmet, then say \emph{hudna} again
\end{itemize}

\begin{tcolorbox}[breakable,colback=teal!4,colframe=teal!35!black,title={Before you move on}]
What does \ar{هدنة} mean? (\textquotedblleft{}A truce\textquotedblright{}.) Say it once more.
\end{tcolorbox}
\section[\texorpdfstring{\ar{معاهدة} (muʿāhada)}{معاهدة (muʿāhada)}]{\ar{معاهدة} (muʿāhada) — a treaty}

\label{lesson:AR-C341-muahada}



\begin{quote}
Before the new one: say the Arabic for a helmet, then the Arabic for a truce.
\end{quote}

\subsection*{You'll want to know: \ar{معاهدة}}
\textbf{\ar{معاهدة}} — \emph{muʿāhada} — \textquotedblleft{}a treaty\textquotedblright{}.

\begin{quote}
\ar{وقع} \ar{البلدان} \ar{معاهدة} \ar{سلام}.
\end{quote}

\begin{quote}
\emph{waqqaʿa al-baladān muʿāhadat salām.}
\end{quote}

\begin{quote}
The two countries signed a peace treaty.
\end{quote}

\begin{cousinweb}
\textbf{\ar{معاهدة}} (\emph{muʿāhada}) is built on the native Arabic root \emph{ʿ-h-d} (Lane), where \emph{ʿahd} is a pledge. \emph{Muʿāhada} is the verbal noun of Form III, \textquotedblleft{}pledging to one another\textquotedblright{}.
\end{cousinweb}

\begin{grammarlens}[title={What you've built}]
One more word for telling what happened.
\end{grammarlens}

\subsection*{Your turn}
\begin{itemize}
\raggedright
  \item \emph{Say it:} \emph{muʿāhada}
  \item \emph{Say it:} \emph{muʿāhada}, once more
  \item \emph{Say it:} \emph{waqqaʿa al-baladān muʿāhadat salām.} — The two countries signed a peace treaty.
  \item \emph{Recall:} say the Arabic for a helmet, then the Arabic for a truce, then say \emph{muʿāhada} again
\end{itemize}

\begin{tcolorbox}[breakable,colback=teal!4,colframe=teal!35!black,title={Before you move on}]
What does \ar{معاهدة} mean? (\textquotedblleft{}A treaty\textquotedblright{}.) Say it once more.
\end{tcolorbox}
\section[\texorpdfstring{\ar{خان} (khāna)}{خان (khāna)}]{\ar{خان} (khāna) — to betray}

\label{lesson:AR-C341-khana}



\begin{quote}
Before the new one: say the Arabic for a truce, then the Arabic for a treaty.
\end{quote}

\subsection*{You'll want to know: \ar{خان}}
\textbf{\ar{خان}} — \emph{khāna} — \textquotedblleft{}to betray\textquotedblright{}.

\begin{quote}
\ar{خان} \ar{الوزير} \ar{الملك} \ar{في} \ar{القصة}.
\end{quote}

\begin{quote}
\emph{khāna al-wazīr al-malik fī al-qiṣṣa.}
\end{quote}

\begin{quote}
In the story the minister betrayed the king.
\end{quote}

Present: \ar{يخون}, \emph{yakhūn}.

Not the noun \ar{خان}, an inn, which has another source.

\begin{cousinweb}
\textbf{\ar{خان}} (\emph{khāna}) is built on the native Arabic root \emph{kh-w-n} (Lane), \textquotedblleft{}to be unfaithful to a trust\textquotedblright{}; \emph{khiyāna} is treachery.
\end{cousinweb}

\begin{grammarlens}[title={What you've built}]
One more word for telling what happened.
\end{grammarlens}

\subsection*{Your turn}
\begin{itemize}
\raggedright
  \item \emph{Say it:} \emph{khāna}
  \item \emph{Say it:} \emph{khāna}, once more
  \item \emph{Say it:} \emph{khāna al-wazīr al-malik fī al-qiṣṣa.} — In the story the minister betrayed the king.
  \item \emph{Recall:} say the Arabic for a truce, then the Arabic for a treaty, then say \emph{khāna} again
\end{itemize}

\begin{tcolorbox}[breakable,colback=teal!4,colframe=teal!35!black,title={Before you move on}]
What does \ar{خان} mean? (\textquotedblleft{}To betray\textquotedblright{}.) Say it once more.
\end{tcolorbox}
